\clearpage
\section{Recomendaciones para modelamiento e implicancias de negocio}

\subsection{Clasificación de variables}

\begin{table}[H]
\centering
\caption{Tratamiento recomendado antes de entrenar.}
\label{tab:tratamiento-modelo}
\begin{tabularx}{\textwidth}{p{3.2cm}p{3.2cm}X}
\toprule
Clasificación & Variables & Recomendación \\
\midrule
Target & \code{popularity} & Regresión; conservar escala 0--100 y analizar masa en cero \\
Recomendadas & Atributos acústicos, \code{explicit} & Línea base y modelos regularizados; escalar cuando corresponda \\
Transformación & \code{duration_ms}, extremos continuos & Comparar original con \code{log1p}; usar pipeline para evitar fuga \\
Encoding & \code{track_genres}, \code{key}, \code{mode}, \code{time_signature} & Multi-hot para géneros; one-hot o codificación categórica para el resto \\
Alta cardinalidad & \code{artists}, \code{album_name}, \code{track_name} & Excluir de línea base; evaluar después con validación estricta y propósito claro \\
Excluir & \code{track_id} & Nunca predictor; usar solo como identificador y control \\
\bottomrule
\end{tabularx}
\end{table}

El pipeline debe ajustar transformaciones solo con entrenamiento. Una línea base que prediga la
mediana permite comprobar si modelos más complejos agregan valor. Después, una regresión lineal o
Ridge ofrece interpretabilidad; alternativas no lineales pueden explorar curvaturas observadas, sin
tuning excesivo en esta fase.

La partición aleatoria es un punto de partida, pero debe complementarse con grupos por artista para
medir generalización a identidades nuevas. Si se obtienen fechas, una partición temporal es más
realista para predecir periodos posteriores. La unicidad actual de \code{track_id} evita una fuga
directa, aunque no la dependencia entre canciones del mismo artista.

MAE se interpreta como error medio en puntos; RMSE revela errores grandes; $R^2$ compara contra la
media y puede ser bajo o negativo fuera de muestra. Deben presentarse intervalos mediante
validación cruzada y métricas por segmentos. Un $R^2$ alto no demuestra causalidad, equidad ni
utilidad operativa.

\subsection{Implicancias de negocio}

El análisis enseña que la popularidad no se explica mediante una regla acústica simple. Algunas
variables contienen señal descriptiva y los géneros segmentan diferencias, pero factores de
exposición probablemente faltan. Una futura estimación podría priorizar revisión humana, comparar
segmentos o identificar casos donde convenga investigar contexto adicional.

No puede afirmarse que cambiar energía, sonoridad, explicitud o instrumentalidad produzca éxito. La
predicción debería acompañarse de información de incertidumbre y nunca determinar por sí sola
presupuesto, lanzamiento o mérito artístico. La decisión útil combina evidencia cuantitativa,
objetivos creativos, conocimiento del mercado y experimentación posterior.
